\chapter{The Five-Step Risk Assessment Method Applied to Construction}
\label{chap:The Five-Step Risk Assessment Method Applied to Construction}

%---------------------------- SECTION ----------------------------
\section{Why this assessment matters in construction}

The five-step risk assessment method is not a bureaucratic formality. It is a structured cognitive process that, when applied with rigour and site-specific knowledge, produces the understanding necessary to keep workers safe and the documentation necessary to demonstrate that the duty of care has been discharged. In construction, where the hazard landscape changes daily and the consequences of uncontrolled risk are frequently fatal, the quality of the risk assessment process is directly and demonstrably linked to incident rates. Sites that produce careful, specific, and regularly reviewed assessments conduct their work within managed risk envelopes; sites that treat risk assessment as a box-ticking exercise to satisfy pre-qualification requirements experience the predictable outcomes of that approach.

The most persistent and damaging failure in construction risk assessment practice is the generic document. Risk assessments produced from templates, downloaded from trade associations, or replicated from previous projects without modification do not constitute suitable and sufficient assessments under Regulation 3 of the Management of Health and Safety at Work Regulations 1999. They may list the correct hazards in the abstract, but they do not reflect the specific site conditions, the specific workforce, the specific method of work, or the specific interface hazards that arise when multiple trades work in proximity. A risk assessment for working at height that does not specify the type of access equipment, the edge protection details, the rescue procedure, and the specific fragile roof areas on the actual project is not a risk assessment; it is a list of hazards with a signature on the bottom.

A second persistent failure is the separation of the risk assessment from the method statement. The risk assessment identifies what can go wrong and assesses the level of risk; the method statement describes how the work will be done in a way that incorporates the controls identified in the assessment. These two documents are logically inseparable, yet they are routinely produced by different people at different times, with the result that the method statement describes a method that does not reflect the controls in the risk assessment, or the risk assessment identifies controls that the method statement does not describe. The two documents must be produced together, reviewed together, and updated together.

A further failure that is particularly prevalent in the construction sector is the inaccessibility of risk assessments to the workforce. Assessments produced in complex technical language and kept in a site office filing cabinet have no protective value for the operative who cannot read English, has not been inducted, and is not aware that the assessment exists. The five-step process must culminate in a communication and briefing step that ensures the key findings --- the significant hazards, the controls in place, and the required behaviours --- are understood by the workers who are exposed to the risks. Toolbox talks, pre-task briefings, and translated safety information are not optional extras; they are the mechanism by which the risk assessment produces its protective effect.

\barquote{``A risk assessment that nobody reads protects nobody. The document is not the goal --- the conversation it generates is.'' --- A recurring lesson from HSE investigations into fatal construction incidents.}

\example{Site Reality}{A demolition contractor produces a risk assessment and method statement (RAMS) for the removal of a suspended ceiling system on the first floor of an office building undergoing refurbishment. The RAMS is a standard template with the project name and address inserted at the top. It lists ``asbestos'' as a hazard with the control noted as ``asbestos survey to be completed''. However, the asbestos management survey that was conducted recorded assumed asbestos-containing materials in the ceiling tiles and the associated textured coating. The refurbishment survey --- which is required before any demolition or intrusive work, and is a different survey type from the management survey --- has not been conducted. The demolition operative, briefed only verbally that ``there might be some asbestos about'', begins removing the suspended ceiling tiles with a claw hammer, releasing fibres from damaged tile edges. The RAMS failed at Step 1 (hazards not specifically identified), Step 4 (controls not implemented --- the required survey was not completed), and Step 5 (not reviewed since project commencement two weeks prior).}

%---------------------------- SECTION ----------------------------
\section{Legal and professional framework}

Regulation 3 of the Management of Health and Safety at Work Regulations 1999 places an explicit duty on every employer and self-employed person to make a suitable and sufficient assessment of the risks to the health and safety of their employees and others who may be affected by their work. The phrase ``suitable and sufficient'' is not defined in the Regulations but has been extensively interpreted by the courts and the HSE: an assessment is suitable and sufficient if it identifies the significant risks arising from the work, identifies the persons at risk, assesses the level of those risks, and records the controls that reduce them to the lowest reasonably practicable level. It must be site-specific and task-specific; it must reflect actual conditions, not hypothetical ones; and it must be produced by a competent person.

HSG65, the HSE's managing for health and safety framework, provides the operational model within which the five-step process sits. The Plan--Do--Check--Act cycle requires that risk assessment is embedded in the planning phase, that controls are implemented in the doing phase, that compliance and effectiveness are checked through monitoring, and that the assessment is reviewed in the act phase when monitoring reveals deficiencies or conditions change. L21, the Approved Code of Practice for the Management Regulations 1999, provides specific guidance on the standard expected of risk assessments in different workplace contexts, including construction. CDM 2015 adds the requirement that risk assessment outputs inform the construction phase plan and that the principal designer ensures pre-construction health and safety information is used in the design and planning of work.

\paragraph{Key frameworks relevant to this chapter include: Management of Health and Safety at Work Regulations 1999 reg.3; HSG65; L21 ACOP; CDM 2015.}

\begin{itemize}
  \bitem{Management of Health and Safety at Work Regulations 1999 reg.3}{Imposes the core duty to conduct suitable and sufficient risk assessments; requires that significant findings are recorded (employers with five or more employees); requires review when the assessment is no longer valid; applies to all employers and self-employed persons regardless of project size.}
  \bitem{HSG65 --- Managing for Health and Safety}{Provides the Plan--Do--Check--Act framework within which risk assessment sits; establishes that risk assessment is a continuous management process, not a one-time event; sets out the principles of proportionate risk management applicable to all construction activities.}
  \bitem{L21 ACOP --- Management of Health and Safety at Work}{The Approved Code of Practice for the Management Regulations, providing specific guidance on what constitutes a suitable and sufficient risk assessment; has quasi-legal status, and departure from its guidance must be justified by equivalent or superior alternative arrangements.}
  \bitem{CDM 2015 --- Construction Phase Plan}{Requires that the construction phase plan produced by the principal contractor includes the site-specific risks identified in the pre-construction information, the arrangements for managing those risks, and the health and safety provisions for the specific project, all of which must be underpinned by site-specific risk assessments.}
\end{itemize}

%---------------------------- SECTION ----------------------------
\section{Conducting the assessment using the five-step method}

\subsection{Step 1 --- Identify hazards}
\paragraph{Hazard identification in construction must be conducted through a combination of desk-based analysis and physical site inspection. The desk-based element draws on the pre-construction information provided under CDM 2015, existing surveys (asbestos, ground investigation, structural), design information, drawings, and specifications. The site inspection element --- often called a hazard walk-down --- must be conducted by a competent person familiar with the specific type of work, ideally accompanied by the supervisor or foreman who will oversee the activity. Hazards must be identified for each distinct phase of the work, for each distinct trade, and for the interfaces between trades. The inspection should consider both obvious physical hazards and the less visible health hazards: RCS from concrete and masonry cutting, asbestos in legacy materials, noise from plant and power tools, vibration from handheld equipment, and chemical exposure from COSHH-regulated substances. The output of Step 1 is a comprehensive hazard register, specific to the site and the activity, which forms the foundation for all subsequent steps.}

\subsection{Step 2 --- Decide who or what is affected}
\paragraph{The assessment must explicitly identify all groups who may be harmed by each hazard identified in Step 1. On a construction site, this population includes: directly employed operatives performing the task; other workers in adjacent areas who may be exposed without performing the task themselves; self-employed persons and subcontractors; supervisors, managers, and designers who visit the working area; delivery drivers; and members of the public, including children, who may be adjacent to the site boundary. Workers with specific vulnerabilities must be individually considered: young workers and apprentices who lack experience; new and expectant mothers; workers with pre-existing health conditions that increase susceptibility to noise, vibration, or dust; and workers who are new to the site or whose first language is not English. The assessment must demonstrate that it has considered each of these groups, not merely that it has identified ``workers'' as a generic category.}

\subsection{Step 3 --- Evaluate likelihood and consequence}
\paragraph{Risk evaluation must record both the inherent risk --- the risk before any controls are applied --- and the residual risk --- the risk remaining after controls are in place. Recording only the residual risk, as is common in many construction risk assessment templates, prevents the assessor and the reader from understanding whether the controls are proportionate to the hazard. A risk matrix should be applied consistently: a five-by-five matrix with likelihood scored 1 (rare) to 5 (almost certain) and consequence scored 1 (minor injury) to 5 (fatality or catastrophic harm) produces a risk rating from 1 to 25, with ratings above 15 typically classified as high risk requiring immediate and substantial control, ratings of 9--15 as medium risk requiring specific controls and monitoring, and ratings below 9 as lower risk manageable through standard site precautions. The residual risk rating must reflect the realistic effectiveness of the controls in place, not the ideal effectiveness of perfectly implemented controls; a control that is frequently bypassed in practice provides less risk reduction than the assessor may assume.}

\subsection{Step 4 --- Select and implement controls}
\paragraph{Controls must be selected by applying the hierarchy of controls in sequence (see Chapter~\ref{chap:What is Risk in a Construction and Trades Context?}), and each control must be specific, actionable, and assigned to a named individual responsible for implementation and ongoing maintenance. Vague controls such as ``ensure safe access'' or ``use appropriate PPE'' do not constitute effective risk controls; they transfer the decision about what to do to the operative at the point of risk, which is precisely the decision-making context in which errors occur. Controls must be written in specific, imperative language: ``erect a fully-boarded scaffold with double guardrails and toe-boards on the north and east elevations before any roof work commences, maintained by [named subcontractor] and inspected weekly''. Each control must be implemented before the activity begins, not after the first incident. The method statement must be updated to reflect the controls selected, so that the operative briefing describes a method that already incorporates the protective measures.}

\subsection{Step 5 --- Record and review}
\paragraph{The written record of a construction risk assessment must include: the project name and location; the specific activity assessed; the date of the assessment and the name of the competent person who conducted it; a list of hazards identified; the persons at risk from each hazard; the inherent risk rating; the controls in place; the residual risk rating; the named owner of each control; the review date; and the signature of the responsible person. The review schedule must be active, not passive: it must be entered into the site management calendar, allocated to a named person, and recorded as completed. Conditions that trigger an immediate unscheduled review include any change to the work method, any new hazard identified on site, any near-miss or incident, any change in the workforce, or any variation to the design or programme that affects the activity assessed. In practice, on a live construction site, this means that the risk assessment review should be a standing item on the weekly site management meeting agenda.}

\example{Five-Step Application}{A groundworks team is tasked with breaking out and removing a reinforced concrete ground slab in preparation for a new basement construction. Step 1 identifies the following hazards: RCS from concrete breaking; noise and vibration from breakers and compactors; manual handling of broken concrete arisings; risk of encountering buried services within the slab or beneath it; and structural risk from undermining adjacent foundations. Step 2 identifies: the groundworkers operating the breakers; the site manager conducting inspections; adjacent trades working in the vicinity; and a structural engineer advising on foundation proximity. Step 3 assesses the inherent risk of RCS as High (score 20), noise as High (score 16), manual handling as Medium (score 9), buried services as High (score 20), and structural risk as Medium (score 12). Step 4 implements: on-tool water suppression for all breaking; half-day maximum breaker use with task rotation; LEV extraction for residual dust; hearing protection zones with hearing protection provided; mechanical handling of arisings using a telehandler; a full utility scan and trial hole programme before breaking commences; and a structural engineer's written confirmation that the breaking sequence is safe relative to adjacent foundations. Step 5 records the assessment, assigns the site manager as review owner, and mandates review before each day's breaking commences via a pre-task briefing.}

%---------------------------- SECTION ----------------------------
\section{Applying the hierarchy of controls}

\paragraph{The five-step risk assessment method is only as effective as the quality of the controls it generates. The hierarchy of controls is the decision-making framework that ensures the assessor works through control options in order of effectiveness, from the most powerful (elimination) to the least powerful (PPE). In construction practice, the temptation to move directly to PPE as a control is strong, because PPE is visible, auditable, and cheap in the short term. The hierarchy demands a more deliberate approach, and the risk assessment documentation must show that higher-order controls have been considered and either implemented or explicitly rejected as not reasonably practicable before lower-order controls are relied upon.}

\begin{itemize}
  \bitem{Elimination}{Can the activity that generates the hazard be removed entirely from the scope? Can the in-situ concrete slab be left in place and the basement design revised? Can the working at height component be designed out by using a ground-level assembly method? Can the asbestos survey be completed before any demolition, removing the uncertainty about fibre release during the work?}
  \bitem{Substitution}{Can the hazardous process be replaced with a less hazardous alternative? Can diamond wire sawing replace percussive breaking to reduce RCS and noise? Can pre-cut kerb stones replace on-site cutting? Can a lower-silica concrete be specified for repairs? Can battery-powered tools with documented lower vibration magnitudes replace pneumatic equivalents?}
  \bitem{Engineering controls}{Can physical measures be implemented to reduce exposure independent of individual behaviour? Can on-tool water suppression be installed on all cutting and breaking equipment? Can local exhaust ventilation capture residual dust? Can acoustic enclosures or barriers be erected around high-noise operations? Can mechanical lifting aids eliminate the manual handling of heavy materials?}
  \bitem{Administrative controls}{Can exposure be limited through planning and procedure? Can task rotation limit daily vibration exposure to below the upper exposure action value? Can permit-to-work systems control access to high-risk areas? Can toolbox talks and pre-task briefings ensure that controls are understood and implemented before work begins? Can health surveillance programmes identify early signs of occupational disease before permanent harm occurs?}
  \bitem{PPE}{As the last resort, what specific PPE is required to control residual risk? FFP3 respirators for RCS; Class 4 hearing protection for noise above the upper exposure action value; anti-vibration gloves to supplement engineering controls for vibration; hard hats, safety footwear, and high-visibility clothing as baseline provisions across all construction activities.}
\end{itemize}

%---------------------------- SECTION ----------------------------
\section{Cadence of assessment and review triggers}

\paragraph{The review cycle for construction risk assessments must reflect the dynamic nature of site conditions. A risk assessment that was suitable and sufficient on day one of a groundworks package may be dangerously inadequate by week three, when the excavation has deepened, adjacent structures have been exposed, and groundwater has been encountered. The minimum review cadence must be established in the construction phase plan and must be actively managed as a site management obligation, not left to the discretion of individual supervisors.}

\begin{itemize}
  \bitem{Pre-task or daily review for high-risk activities}{Activities with a residual risk rating of high or very high --- working at height, confined space entry, excavation below three metres, hot work, live electrical work --- must be reviewed before each day's work commences via a pre-task briefing that confirms conditions match those described in the risk assessment.}
  \bitem{Weekly review for ongoing medium-risk activities}{Activities with a medium residual risk rating that are ongoing throughout a project phase should be reviewed weekly, as part of the regular site inspection programme, to confirm that controls remain in place and effective.}
  \bitem{Immediate review on any change to method, plant, or materials}{Any change to the work method, the plant or equipment used, the materials involved, or the personnel performing the activity must trigger an immediate review and update of the assessment before the changed activity commences.}
  \bitem{Immediate review following any near-miss or incident}{Any near-miss reported through the site safety management system, or any incident resulting in first-aid treatment or above, must trigger an immediate review of the risk assessment applicable to the activity involved.}
  \bitem{Review at each project phase change}{The transition from one major phase to the next (groundworks to structure; structure to envelope; envelope to fit-out) must trigger a comprehensive review of all current risk assessments, as the hazard profile of the site changes fundamentally at each phase.}
\end{itemize}

%---------------------------- SECTION ----------------------------
\section{Common failure modes}

\begin{itemize}
  \bitem{Generic and unmodified template documents}{The most common failure: a template downloaded, minimally edited, and submitted as a site-specific assessment. The template does not reflect actual site conditions, actual hazard interactions, or actual controls. It will not identify interface hazards between trades, site-specific ground conditions, or the specific equipment being used.}
  \bitem{No inherent risk recorded}{Assessments that record only residual risk make it impossible to judge whether controls are proportionate to the hazard. A residual risk of ``low'' following the implementation of PPE as the sole control for an activity with a very high inherent risk is not a well-controlled risk; it is an incorrectly scored risk.}
  \bitem{Controls without named owners or timescales}{Controls listed without a named individual responsible for implementation and a date by which they will be in place are aspirational, not operational. On a multi-contractor site, unnamed controls are invariably assumed by each contractor to be someone else's responsibility.}
  \bitem{Language inaccessible to the workforce}{Risk assessments written in dense professional language and not translated or communicated in a format accessible to operatives whose first language is not English do not fulfil the purpose of the assessment. The operative at risk cannot act on information they cannot access or understand.}
  \bitem{Method statement and risk assessment misaligned}{Method statements that describe methods inconsistent with the controls in the associated risk assessment expose the workforce to uncontrolled risks during normal working. The two documents must be produced together, reviewed together, and updated together.}
  \bitem{Review cycle not implemented or not documented}{Risk assessments with passed review dates are a recurring finding in post-incident HSE investigations. The review must actually happen, be documented, and be retained as evidence. A review that took place but was not recorded is legally indistinguishable from a review that did not take place.}
\end{itemize}

\example{Failure Pattern}{A roofing subcontractor submits a risk assessment and method statement for the installation of a standing seam metal roof on a commercial building. The RAMS is dated from a project completed eighteen months previously; the site address has been updated but the hazard identification has not been revised. The assessment does not reference the fragile rooflights present on this building, does not identify the 7.5-metre unprotected eave on the west elevation, and specifies ``harness and lanyard'' as the primary fall protection control without identifying suitable anchor points or a rescue procedure. The method statement describes an access method using mobile elevating work platforms, but the site conditions (soft ground, site slope) mean that MEWPs cannot be used and the roofer has instead elected to use a ladder. The mismatch between the method statement, the risk assessment, and the actual method of work means that no effective fall protection controls are in place when an operative falls from the unprotected eave.}

%---------------------------- CHAPTER SUMMARY ----------------------------
\section{Chapter Summary}

\begin{table}[H]
\centering
\begin{tabularx}{\textwidth}{>{\raggedright\arraybackslash}p{3.2cm}X}
\toprule
\textbf{Assessment Element} & \textbf{Operational Requirement} \\
\midrule
Legal/Professional Basis & Management Regulations 1999 reg.3; L21 ACOP; HSG65; CDM 2015. The duty to conduct a suitable and sufficient assessment applies to every employer and self-employed person on every project. \\
Method & Apply the full five-step process and document inherent/residual risk. \\
Controls & Specific, actionable controls with named owners; aligned with the method statement; applied in hierarchy of controls sequence; verified as implemented before work commences. \\
Cadence & Daily pre-task briefing for high-risk activities; weekly review for ongoing medium-risk activities; immediate review on change, near-miss, or incident; comprehensive review at every phase change. \\
Assurance & Use audit, incident learning, and governance reporting to verify effectiveness. \\
\bottomrule
\end{tabularx}
\end{table}

\begin{itemize}
  \bitem{Key takeaway 1}{A suitable and sufficient risk assessment under the Management Regulations 1999 must be site-specific, task-specific, and phase-specific; generic template documents are not compliant and will not withstand scrutiny following an incident.}
  \bitem{Key takeaway 2}{Both inherent and residual risk must be recorded; recording only residual risk conceals whether controls are proportionate to the hazard and makes the assessment unintelligible as a management tool.}
  \bitem{Key takeaway 3}{The risk assessment and method statement are logically inseparable documents that must be produced together, reviewed together, and updated together; misalignment between the two is a systemic failure that leaves workers exposed to uncontrolled risks.}
  \bitem{Key takeaway 4}{The protective value of a risk assessment is only realised when its contents are communicated to and understood by the workers it is intended to protect; inaccessible documents in site office filing cabinets do not protect anyone.}
\end{itemize}

\barquote{In Chapter~\ref{chap:Roles, Responsibilities and Governance in Construction}, we turn from this assessment to the next domain in the Prudentia framework.}
